\section{The DepthShift benchmark}
\label{sec:benchmark}

DepthShift pairs tasks that have an exact per-instance depth oracle with a split design that
separates \emph{longer} inputs from \emph{deeper} ones (\cref{fig:splits} in \cref{app:results}). All generators run on the GPU
from fixed seeds. Test sets are frozen and stored with SHA-256 hashes, and the training sampler
discards any instance whose hash matches an in-distribution validation or test instance.

\paragraph{T1: addition.} Each instance consists of two $L$-digit numbers $a$ and $b$, written
least-significant digit first and encoded as one token $10a_i+b_i$ per position, plus one extra position
for the final carry. By \cref{prop:oracles}, the depth is one plus the longest run of digit pairs that
sum to~9, because a carry generated at a stop must ripple through the whole run before the positions
behind it are determined. Uniformly random digits make long runs exponentially rare (each pair sums
to~9 with probability $1/10$), so we control $\dstar$ directly. The generator inserts exactly one run of
length $\dstar-1$ between two stops, draws background runs at a random density, and caps them below
$\dstar-1$. Stops are drawn uniformly from the 90 pairs with $a+b\ne 9$, which means a run receives an
incoming carry with probability $1/2$.

\paragraph{T2: segmented parity.} The input is a bit string interleaved with delimiter tokens, and the
target at each position is the parity of the current segment up to that position. The depth is the
length of the longest segment. Parity differs from addition in a way that turns out to matter: as
information arrives from the left, the intermediate predictions keep flipping, so the per-step outputs
do not settle early.

\paragraph{T3: mazes.} We also generate $n\times n$ mazes with a randomized depth-first search; the
input channels mark open cells, the start and the goal, and the target is the shortest path. We
record the optimal path length and the number of dead-end-filling iterations a local rule needs to
isolate the path as depth proxies. The generator and splits are released, but within our compute
budget no maze backbone passed Gate G1 (\cref{sec:limitations}), so all experiments below use T1 and
T2.

\paragraph{Splits.} \Cref{tab:splits} summarises the splits for T1 and T2. In-distribution (ID)
instances have lengths 4--20 and depths 1--10. OOD-depth keeps these lengths and raises the depth to
11--20; OOD-length raises the length to 64--256 at ID depths; OOD-both combines lengths 64--256 with
depths 32--128, up to $12.8\times$ the training maximum. A separate OOD-validation split (lengths
24--32, depths 12--16) is used only to tune the backbone, and the thresholds of every stopping rule are
chosen on ID-validation. For mazes we train on $9\times 9$ grids and test on $13\times13$ up to
$33\times33$.

\begin{table}[t]
  \centering
  \caption{DepthShift splits for addition (T1) and segmented parity (T2). Every OOD split is evaluated
  once, at the end; thresholds are selected on ID-val.}
  \label{tab:splits}
  \small
  \begin{tabular}{lrrrl}
    \toprule
    Split & Length & Depth $\dstar$ & Instances & Purpose \\
    \midrule
    Train (sampled online) & 4--20 & 1--10 & -- & training \\
    ID-val / ID-test & 4--20 & 1--10 & 2000 / 2000 & thresholds / test \\
    OOD-val & 24--32 & 12--16 & 1000 & backbone tuning only \\
    OOD-depth & 4--20\textsuperscript{a} & 11--20 & 2000 & deeper, same length \\
    OOD-length & 64--256 & 1--10 & 1000 & longer, same depth \\
    OOD-both & 64--256 & 32--128 & 1000 & deeper and longer \\
    \bottomrule
  \end{tabular}\\[2pt]
  {\footnotesize \textsuperscript{a}10--20 for addition and 11--20 for parity, the shortest lengths that admit the depth.}
\end{table}
